
\chapter{The \texttt{qua} Command-Line Tool}

The \texttt{qua} command-line tool
	is the primary and recommended method for executing Quarb queries
	against structured data files.
It is a powerful, structure-aware replacement
	for traditional text-based search tools like \texttt{grep},
	designed for the modern era of JSON, HTML, XML, and other tree-dominated data formats.


\section{Overview}

\texttt{qua} is a lightweight wrapper around the \texttt{quarb} engine.
It translates your Quarb query string,
	applies it to the parsed structure of an input file,
	and displays the results intuitively.

\paragraph{Key Features:}
\begin{itemize}
	\item \emph{Structure-Aware Search}:
		Understands the hierarchy and relationships in your data,
			eliminating the false positives
			and complex regex patterns required by line-based tools.

	\item \emph{Rich Output}:
		Returns typed results
			(strings, numbers, booleans) instead of just lines of text.

	\item \emph{Powerful Pipelines}:
		Seamlessly integrates into Unix shell pipelines,
			working alongside tools like
			\texttt{jq}, \texttt{sort}, and \texttt{uniq}.
\end{itemize}

\section{Installation}

Install \texttt{qua} using the Rust package manager Cargo:
\begin{verbatim}
cargo install qua
\end{verbatim}
\emph{Other installation methods (e.g., package managers like Homebrew)
	are planned in the future.
	Check the official \texttt{qua} manual for the latest instructions.}

\section{Basic Usage}

The basic command structure is:
\begin{verbatim}
qua [OPTIONS] <QUERY_STRING> [FILE]
\end{verbatim}
If no \texttt{[FILE]} is provided, \texttt{qua} reads from standard input (stdin), allowing it to be used in pipelines.

\section{Examples}
\begin{enumerate}
	\item Find all \texttt{name} fields of users in London from a JSON file:
	\begin{verbatim}
	qua '^//users/*[::city="London"]::name' data.json
	\end{verbatim}

	\item Find the text content of all heading elements in an HTML document from a URL:
	\begin{verbatim}
    curl -s https://example.com | qua '//h1::text'
	\end{verbatim}

	\item Count the number of deprecated functions in a codebase (using a hypothetical AST format):
	\begin{verbatim}
    qua '//function[::is_deprecated]' src/ | wc -l
	\end{verbatim}
\end{enumerate}

\section{Learning More}

This specification defines the Quarb language that \texttt{qua} executes.
For complete documentation on the \texttt{qua} tool itself,
	including command-line options, supported formats,
	and advanced usage,
	please see the official \emph{\texttt{qua} Manual}.

\begin{itemize}
	\item \texttt{qua} Manual: [\texttt{https://github.com/quarb/qua}](https://github.com/quarb/qua)
	\item Quick Start Guide: [\texttt{https://quarb.org/guides/qua}](https://quarb.org/guides/qua)
\end{itemize}
